\chapter{Subarachnoid hemorrhage}
\index{Subarachnoid hemorrhage}

\medskip
\noindent\textit{Educational reference; always seek professional advice for individual care.}

\section*{Definition and Overview}
A subarachnoid hemorrhage (SAH) is bleeding into the subarachnoid space, the fluid-filled space between the brain and the thin membranes that cover it, where the cerebrospinal fluid circulates. It is a form of stroke, accounting for roughly one in twenty strokes, and it is one of the most serious conditions in medicine. The classic picture is a sudden, explosive headache, often described as "the worst headache of my life" or "like being hit on the head", that reaches maximum intensity within seconds to a minute. Most cases, about 80 percent, are caused by the rupture of a brain aneurysm, a balloon-like bulge in the wall of an artery that weakens and bursts. The hemorrhage floods the space around the brain, raising pressure and reducing blood flow, and can cause lasting brain damage or death. In many countries, subarachnoid hemorrhage is a medical emergency: people with suspected SAH are transported rapidly to hospital, because prompt diagnosis and treatment greatly improve the chances of survival and recovery. This chapter explains what happens in the brain, who is at risk, the symptoms that must never be ignored, how the condition is diagnosed and treated, and the road to recovery.

\section*{Epidemiology}
Subarachnoid hemorrhage is relatively uncommon but devastating. Worldwide it affects roughly 6 to 9 people per 100,000 per year, which in many countries means several hundred to over a thousand cases annually. It differs from the commoner forms of stroke in striking younger people: the average age of onset is around 50 to 60, and it can occur in people in their twenties and thirties who are otherwise healthy. It is slightly more common in women than men, and the incidence rises with age until around 60, before declining. Rates vary with ethnicity and geography, and the population includes groups with higher risk. Smoking, high blood pressure, and excessive alcohol use are important modifiable risk factors, and these help explain trends over time: as smoking rates have fallen in many countries, some studies suggest the incidence of SAH has gradually declined. Approximately 10 to 15 percent of people die before reaching hospital, and of those admitted, a significant proportion die or are left disabled, making SAH one of the most feared neurological emergencies.

\section*{Aetiology and Pathophysiology}
The immediate event is the rupture of an artery on the surface of the brain, pouring blood into the subarachnoid space. In most cases the source is a cerebral aneurysm: a localised weakness in the arterial wall, often at branch points, that balloons outward over years. Not all aneurysms rupture; many are found incidentally and never bleed. When rupture does occur, the blood raises intracranial pressure, irritates the brain and its blood vessels, and can trigger spasm of arteries (vasospasm) days later, which further starves the brain of blood.

\begin{itemize}
    \item \textbf{Ruptured cerebral aneurysm:} the cause of about 80 percent of spontaneous SAH; most commonly at the circle of Willis at the base of the brain
    \item \textbf{Hypertension:} high blood pressure weakens artery walls and is a major treatable risk factor
    \item \textbf{Smoking:} damages and weakens blood vessel walls and strongly increases rupture risk
    \item \textbf{Alcohol:} heavy drinking raises blood pressure and increases risk
    \item \textbf{Genetic and connective tissue disorders:} polycystic kidney disease, Ehlers-Danlos syndrome, Marfan syndrome, and a family history of aneurysm all increase risk
    \item \textbf{Vascular malformations:} arteriovenous malformations can bleed into the subarachnoid space, especially in younger people
    \item \textbf{Head injury:} trauma can tear surface vessels and cause SAH without an aneurysm
    \item \textbf{Drug use:} stimulants such as cocaine and methamphetamine can spike blood pressure and trigger rupture
    \item \textbf{Blood thinners:} anticoagulant and antiplatelet medications can worsen bleeding if rupture occurs
\end{itemize}

\section*{Clinical Features}
The symptoms of SAH are dramatic and usually unmistakable, but recognising them quickly saves lives.

\begin{itemize}
    \item \textbf{Sudden "thunderclap" headache:} the hallmark symptom, reaching maximal severity within a second to a minute, unlike any previous headache
    \item \textbf{Nausea and vomiting:} very common in the early phase
    \item \textbf{Loss of consciousness:} about half of patients lose consciousness briefly at the onset, sometimes with a seizure
    \item \textbf{Neck stiffness:} develops over hours as blood irritates the meninges and nerve roots
    \item \textbf{Light sensitivity:} photophobia from meningeal irritation
    \item \textbf{Altered mental state:} confusion, drowsiness, or agitation
    \item \textbf{Neurological deficits:} weakness, numbness, double vision, or difficulty speaking, depending on the location of the bleed
    \item \textbf{Eye changes:} bleeding behind the eye (retinal or subhyaloid hemorrhage), which doctors can detect with an ophthalmoscope
    \item \textbf{Seizures:} occur in a minority and are more likely with surface bleeding
    \item \textbf{Sentinel headache:} a smaller, transient headache days or weeks before the rupture, caused by a "warning leak", which can be the only clue that a catastrophe is coming
\end{itemize}

\section*{Differential Diagnosis}
Several conditions cause a sudden severe headache and must be distinguished from SAH.

\begin{itemize}
    \item \textbf{Migraine:} the commonest mimic, but usually builds gradually over minutes to hours, is throbbing, and typically recurs; thunderclap onset strongly favours SAH
    \item \textbf{Tension-type headache:} gradual, mild-to-moderate, band-like; never reaches maximal intensity in seconds
    \item \textbf{Meningitis:} headache, fever, and neck stiffness, but with a more gradual onset; distinguished by lumbar puncture and blood tests
    \item \textbf{Cervicogenic headache and neck pain:} sudden neck problems can radiate to the head but lack the explosive quality
    \item \textbf{Reversible cerebral vasoconstriction syndrome:} a rare cause of thunderclap headaches with normal imaging apart from narrowed vessels
    \item \textbf{Pituitary apoplexy:} sudden bleeding into a pituitary tumour, causing headache and visual disturbance
    \item \textbf{Cervical artery dissection:} a tear in a neck artery causing sudden headache and stroke symptoms, an important mimic in younger people
\end{itemize}

\section*{When to Seek Medical Care}
A sudden, severe headache reaching maximum intensity in under a minute is a medical emergency, even if it fades. Anyone experiencing this, or a headache with vomiting, neck stiffness, collapse, or neurological symptoms, needs immediate assessment in an emergency department. People with a known unruptured aneurysm, or a family history of aneurysm, should be managed by a specialist team, and any new severe headache in them should be taken very seriously. A sentinel "warning leak" headache can precede a major rupture by days, so even a headache that resolved but was unlike anything before warrants prompt review. In many countries, the ambulance service should be called for anyone suspected of having a subarachnoid hemorrhage, because rapid transport and early treatment improve outcomes.

\textbf{Contact your local emergency services for:}
\begin{itemize}
    \item A sudden, explosive headache, the worst of your life, reaching maximum intensity in seconds
    \item Loss of consciousness, collapse, or a seizure at the time of a sudden headache
    \item Headache with neck stiffness, vomiting, confusion, weakness, or double vision
    \item Any sudden neurological symptoms, which may indicate stroke
\end{itemize}

\section*{Diagnosis and Investigations}
Diagnosis is urgent and definitive, using brain imaging and, if needed, examination of the spinal fluid.

\begin{itemize}
    \item \textbf{CT scan of the head:} performed immediately; detects blood in the subarachnoid space in over 95 percent of cases if done within hours of rupture
    \item \textbf{CT angiography:} a contrast study of the brain's arteries, done at the same time to locate the aneurysm
    \item \textbf{Lumbar puncture:} if the CT scan is normal but SAH is strongly suspected, cerebrospinal fluid is examined for blood and its breakdown products; this is time-sensitive
    \item \textbf{Cerebral angiography:} the most detailed view of the blood vessels, used to plan treatment and find the aneurysm
    \item \textbf{MRI and MR angiography:} used in selected cases, particularly for smaller bleeds or when CT is inconclusive
    \item \textbf{ECG and cardiac monitoring:} heart rhythm disturbances are common after SAH and are monitored
    \item \textbf{Neurological monitoring:} repeated examinations assess conscious state, pupils, and limb function, and detect vasospasm or rebleeding
    \item \textbf{Screening:} relatives of people with familial aneurysms may be offered screening to detect unruptured aneurysms early
\end{itemize}

\section*{Management}
Treatment has two goals: stop further bleeding by securing the aneurysm, and protect the brain during the critical early weeks.

\begin{itemize}
    \item \textbf{Immediate stabilisation:} blood pressure control, pain relief, fluids, and careful monitoring in a specialised stroke or neurosurgical unit
    \item \textbf{Surgical clipping:} an operation places a small metal clip across the neck of the aneurysm, sealing it off from the circulation
    \item \textbf{Endovascular coiling:} a less invasive procedure in which catheters are threaded from a leg artery to the brain and platinum coils are packed into the aneurysm to block it
    \item \textbf{Medications:} drugs to lower blood pressure, prevent vasospasm (such as nimodipine), and manage pain and seizures
    \item \textbf{Treatment of vasospasm:} the artery spasm that peaks around days four to fourteen is treated with careful blood pressure support, medication, and sometimes angioplasty
    \item \textbf{Draining hydrocephalus:} excess spinal fluid that cannot drain normally may be diverted with an external drain
    \item \textbf{Rehabilitation:} physiotherapy, occupational therapy, and speech therapy begin early and continue through recovery
    \item \textbf{Long-term follow-up:} survivors are followed with imaging to check the treated aneurysm and other vessels, and are advised on blood pressure control and lifestyle
\end{itemize}

\section*{Complications}
The complications of SAH occur mostly in the first two weeks and determine much of the outcome.

\begin{itemize}
    \item \textbf{Rebleeding:} an untreated aneurysm can bleed again within hours or days, carrying a very high death rate; securing the aneurysm early prevents this
    \item \textbf{Vasospasm and delayed cerebral ischaemia:} artery spasm weeks after the bleed can cause further strokes; it is the leading cause of death and disability in survivors
    \item \textbf{Hydrocephalus:} blood blocks the flow of cerebrospinal fluid, raising pressure and causing confusion and drowsiness
    \item \textbf{Elevated intracranial pressure:} swelling and blood mass can dangerously compress the brain
    \item \textbf{Seizures:} occur in a minority and are managed with medication
    \item \textbf{Medical complications:} pneumonia, deep vein thrombosis, infections, and heart rhythm disturbances complicate the intensive care phase
    \item \textbf{Neuropsychological effects:} cognitive problems, mood changes, anxiety, depression, and fatigue are common and often the most disabling long-term consequences
\end{itemize}

\section*{Prognosis and Outcomes}
Subarachnoid hemorrhage is a grave condition, but outcomes have improved steadily with modern emergency and neurosurgical care. Roughly 10 to 15 percent of people die before reaching hospital, and of those who reach hospital, around a quarter die and a similar proportion are left with significant disability. However, for people who reach hospital alive and have their aneurysm secured quickly, the outlook is far better than it was decades ago, and a substantial proportion, perhaps a third to a half of survivors, make a good functional recovery and return to independent life. Recovery is often slow and can continue for a year or more. Early survivors commonly struggle with fatigue, concentration, memory, and mood, and rehabilitation and psychological support are central to recovery. Age, the amount of bleeding, conscious state on arrival, and the development of vasospasm are the main predictors of outcome. Among people found to have an unruptured aneurysm, the annual rupture risk is low for small aneurysms, and careful specialist discussion weighs treatment against observation. With risk-factor control, particularly stopping smoking and treating high blood pressure, the chance of a first or repeat hemorrhage falls substantially.

\section*{Prevention and Self-Care}
\begin{itemize}
    \item Control high blood pressure with regular checks, medication, exercise, and a low-salt diet
    \item Do not smoke; stopping smoking is one of the most effective ways to reduce aneurysm risk
    \item Limit alcohol to within national guidelines, and avoid binge drinking
    \item Avoid stimulant drugs, which can spike blood pressure and trigger rupture
    \item Seek urgent medical help for any sudden severe headache, including one that resolves, since it may be a warning leak
    \item If an aneurysm is found incidentally, follow the specialist plan for monitoring or treatment and keep blood pressure tightly controlled
    \item Know your family history and discuss screening with a doctor if close relatives have had aneurysms
    \item Support recovery after a hemorrhage with rehabilitation, gradual activity, and attention to mood and cognition
\end{itemize}

\section*{Sources and Further Reading}
The following organisations provide reliable information about subarachnoid hemorrhage, stroke, and brain aneurysms.

\begin{itemize}
    \item healthdirect Australia. \textit{Subarachnoid haemorrhage and brain aneurysm.} https://www.healthdirect.gov.au/
    \item Stroke Foundation (Australia). \textit{Stroke types and information for patients and carers.} https://strokefoundation.org.au/
    \item NHS. \textit{Subarachnoid haemorrhage: symptoms and treatment.} https://www.nhs.uk/
    \item Mayo Clinic. \textit{Brain aneurysm and subarachnoid hemorrhage.} https://www.mayoclinic.org/
    \item World Health Organization (WHO). \textit{Cerebrovascular disease and stroke fact sheets.} https://www.who.int/
    \item Australian Department of Health and Aged Care. \textit{Stroke and cardiovascular disease resources.} https://www.health.gov.au/
\end{itemize}

\clearpage
